% preamble.tex -- packages, theorem environments and macros shared by every chapter.
% Loaded by AnalysisBook.tex with % preamble.tex -- packages, theorem environments and macros shared by every chapter.
% Loaded by AnalysisBook.tex with % preamble.tex -- packages, theorem environments and macros shared by every chapter.
% Loaded by AnalysisBook.tex with % preamble.tex -- packages, theorem environments and macros shared by every chapter.
% Loaded by AnalysisBook.tex with \input{preamble}.

\hypersetup{colorlinks}

% ---------------------------------------------------------------- Packages
% (The tufte-book class already loads hyperref, xcolor, paralist, titlesec,
% titletoc, multicol and natbib.)
\usepackage{amsmath, amsthm, amssymb}
\usepackage{mathrsfs}
\usepackage{booktabs}
\usepackage{enumerate}
\usepackage{manfnt}   % the "dangerous bend" sign used by the Caution environment
\usepackage{graphicx}
\setkeys{Gin}{width=\linewidth,totalheight=\textheight,keepaspectratio}
\graphicspath{{graphics/}}

\usepackage{makeidx}
\makeindex

% Cross-references to figures and tables.  Write \caption{...}\label{...} as
% usual.  The Tufte class sets a float's caption only at \end{figure}, and
% postpones the \label to match; but amsmath resets \label at the start of
% every minipage, which silently undoes that, so \ref would print the current
% section or theorem number instead of the figure number.  This restores it.
\makeatletter
\expandafter\apptocmd\csname\string\@tufte@float\endcsname
  {\let\label\@tufte@label}{}{\typeout{Warning: tufte float label patch failed}}
\makeatother

% ---------------------------------------------------------------- Figures (TikZ)
% Every diagram in the book is drawn with these shared styles, so that the
% pictures look alike and one convention holds throughout:
%   a FILLED endpoint belongs to the set, a HOLLOW endpoint does not.
\usepackage{tikz}
\usetikzlibrary{arrows.meta,calc,decorations.pathreplacing}
\definecolor{accent}{HTML}{A63D2F}% brick red: the set under discussion
\definecolor{calm}{HTML}{2B6A8C}%   slate blue: neighborhoods and covers
\tikzset{
  >={Stealth[length=4pt,width=3.5pt]},
  numline/.style={draw=black!55, line width=0.5pt, <->},
  tickmark/.style={draw=black!55, line width=0.5pt},
  setline/.style={draw=accent, line width=1.8pt},
  nbhdline/.style={draw=calm, line width=1.8pt},
  closedpt/.style={circle, fill=#1, draw=#1, line width=0.7pt, inner sep=0pt, minimum size=4pt},
  closedpt/.default=accent,
  openpt/.style={circle, fill=white, draw=#1, line width=0.7pt, inner sep=0pt, minimum size=4pt},
  openpt/.default=accent,
  seqpt/.style={circle, fill=black, inner sep=0pt, minimum size=2.6pt},
  lbl/.style={font=\footnotesize, inner sep=1.5pt},
  bracedeco/.style={decorate, decoration={brace, amplitude=3.5pt}, line width=0.4pt, draw=black!70},
  % graphs of functions
  graphaxis/.style={draw=black!55, line width=0.5pt, ->},
  curve/.style={draw=accent, line width=1.1pt},
  guide/.style={draw=black!45, dashed, line width=0.4pt},
  epsband/.style={fill=calm!16, draw=none},%   a tolerance on the output
  delband/.style={fill=accent!14, draw=none},% a tolerance on the input
}

% ---------------------------------------------------------------- Numbering
\setcounter{secnumdepth}{2}
\renewcommand{\thesection}{\thechapter.\arabic{section}}
\numberwithin{section}{chapter}

% \bookchapter{One}{Title} starts "Chapter One: Title".  Stepping the chapter
% counter here resets the section, theorem, definition, proposition, figure
% and equation counters automatically, so chapters need no manual \setcounter.
\newcommand{\bookchapter}[2]{%
  \stepcounter{chapter}%
  \chapter{Chapter #1: #2}}

% ---------------------------------------------------------------- Theorem environments
\newtheorem{theorem}{Theorem}[chapter]
\newtheorem{lemma}[theorem]{Lemma}
\newtheorem{corollary}[theorem]{Corollary}
\newtheorem{proposition}{Proposition}[chapter]
\theoremstyle{definition}
\newtheorem{definition}{Definition}[chapter]
\newtheorem*{remark}{Remark}
\newtheorem*{example}{Example}
\newtheorem*{imp}{Important point}
\newtheorem*{dangerousbend}{\textdbend Caution}
\newtheorem*{summary}{Summary}

% Original end-of-chapter problems: \begin{problem}[Optional title] ... \end{problem}
\newtheorem{problem}{Problem}[chapter]
% A hint, set in the margin beside a problem.
% Put it straight after \begin{problem}[...] so it lines up with the heading.
\newcommand{\hint}[1]{\leavevmode\marginnote{\emph{Hint.} #1}\ignorespaces}

% ---------------------------------------------------------------- Front-matter macros
% A full-width epigraph: \openepigraph{quotation}{author}
\newcommand{\openepigraph}[2]{%
  \begin{fullwidth}
  \sffamily\large
  \begin{doublespace}
  \noindent\allcaps{#1}\\% epigraph
  \noindent\allcaps{#2}% author
  \end{doublespace}
  \end{fullwidth}
}
.

\hypersetup{colorlinks}

% ---------------------------------------------------------------- Packages
% (The tufte-book class already loads hyperref, xcolor, paralist, titlesec,
% titletoc, multicol and natbib.)
\usepackage{amsmath, amsthm, amssymb}
\usepackage{mathrsfs}
\usepackage{booktabs}
\usepackage{enumerate}
\usepackage{manfnt}   % the "dangerous bend" sign used by the Caution environment
\usepackage{graphicx}
\setkeys{Gin}{width=\linewidth,totalheight=\textheight,keepaspectratio}
\graphicspath{{graphics/}}

\usepackage{makeidx}
\makeindex

% Cross-references to figures and tables.  Write \caption{...}\label{...} as
% usual.  The Tufte class sets a float's caption only at \end{figure}, and
% postpones the \label to match; but amsmath resets \label at the start of
% every minipage, which silently undoes that, so \ref would print the current
% section or theorem number instead of the figure number.  This restores it.
\makeatletter
\expandafter\apptocmd\csname\string\@tufte@float\endcsname
  {\let\label\@tufte@label}{}{\typeout{Warning: tufte float label patch failed}}
\makeatother

% ---------------------------------------------------------------- Figures (TikZ)
% Every diagram in the book is drawn with these shared styles, so that the
% pictures look alike and one convention holds throughout:
%   a FILLED endpoint belongs to the set, a HOLLOW endpoint does not.
\usepackage{tikz}
\usetikzlibrary{arrows.meta,calc,decorations.pathreplacing}
\definecolor{accent}{HTML}{A63D2F}% brick red: the set under discussion
\definecolor{calm}{HTML}{2B6A8C}%   slate blue: neighborhoods and covers
\tikzset{
  >={Stealth[length=4pt,width=3.5pt]},
  numline/.style={draw=black!55, line width=0.5pt, <->},
  tickmark/.style={draw=black!55, line width=0.5pt},
  setline/.style={draw=accent, line width=1.8pt},
  nbhdline/.style={draw=calm, line width=1.8pt},
  closedpt/.style={circle, fill=#1, draw=#1, line width=0.7pt, inner sep=0pt, minimum size=4pt},
  closedpt/.default=accent,
  openpt/.style={circle, fill=white, draw=#1, line width=0.7pt, inner sep=0pt, minimum size=4pt},
  openpt/.default=accent,
  seqpt/.style={circle, fill=black, inner sep=0pt, minimum size=2.6pt},
  lbl/.style={font=\footnotesize, inner sep=1.5pt},
  bracedeco/.style={decorate, decoration={brace, amplitude=3.5pt}, line width=0.4pt, draw=black!70},
  % graphs of functions
  graphaxis/.style={draw=black!55, line width=0.5pt, ->},
  curve/.style={draw=accent, line width=1.1pt},
  guide/.style={draw=black!45, dashed, line width=0.4pt},
  epsband/.style={fill=calm!16, draw=none},%   a tolerance on the output
  delband/.style={fill=accent!14, draw=none},% a tolerance on the input
}

% ---------------------------------------------------------------- Numbering
\setcounter{secnumdepth}{2}
\renewcommand{\thesection}{\thechapter.\arabic{section}}
\numberwithin{section}{chapter}

% \bookchapter{One}{Title} starts "Chapter One: Title".  Stepping the chapter
% counter here resets the section, theorem, definition, proposition, figure
% and equation counters automatically, so chapters need no manual \setcounter.
\newcommand{\bookchapter}[2]{%
  \stepcounter{chapter}%
  \chapter{Chapter #1: #2}}

% ---------------------------------------------------------------- Theorem environments
\newtheorem{theorem}{Theorem}[chapter]
\newtheorem{lemma}[theorem]{Lemma}
\newtheorem{corollary}[theorem]{Corollary}
\newtheorem{proposition}{Proposition}[chapter]
\theoremstyle{definition}
\newtheorem{definition}{Definition}[chapter]
\newtheorem*{remark}{Remark}
\newtheorem*{example}{Example}
\newtheorem*{imp}{Important point}
\newtheorem*{dangerousbend}{\textdbend Caution}
\newtheorem*{summary}{Summary}

% Original end-of-chapter problems: \begin{problem}[Optional title] ... \end{problem}
\newtheorem{problem}{Problem}[chapter]
% A hint, set in the margin beside a problem.
% Put it straight after \begin{problem}[...] so it lines up with the heading.
\newcommand{\hint}[1]{\leavevmode\marginnote{\emph{Hint.} #1}\ignorespaces}

% ---------------------------------------------------------------- Front-matter macros
% A full-width epigraph: \openepigraph{quotation}{author}
\newcommand{\openepigraph}[2]{%
  \begin{fullwidth}
  \sffamily\large
  \begin{doublespace}
  \noindent\allcaps{#1}\\% epigraph
  \noindent\allcaps{#2}% author
  \end{doublespace}
  \end{fullwidth}
}
.

\hypersetup{colorlinks}

% ---------------------------------------------------------------- Packages
% (The tufte-book class already loads hyperref, xcolor, paralist, titlesec,
% titletoc, multicol and natbib.)
\usepackage{amsmath, amsthm, amssymb}
\usepackage{mathrsfs}
\usepackage{booktabs}
\usepackage{enumerate}
\usepackage{manfnt}   % the "dangerous bend" sign used by the Caution environment
\usepackage{graphicx}
\setkeys{Gin}{width=\linewidth,totalheight=\textheight,keepaspectratio}
\graphicspath{{graphics/}}

\usepackage{makeidx}
\makeindex

% Cross-references to figures and tables.  Write \caption{...}\label{...} as
% usual.  The Tufte class sets a float's caption only at \end{figure}, and
% postpones the \label to match; but amsmath resets \label at the start of
% every minipage, which silently undoes that, so \ref would print the current
% section or theorem number instead of the figure number.  This restores it.
\makeatletter
\expandafter\apptocmd\csname\string\@tufte@float\endcsname
  {\let\label\@tufte@label}{}{\typeout{Warning: tufte float label patch failed}}
\makeatother

% ---------------------------------------------------------------- Figures (TikZ)
% Every diagram in the book is drawn with these shared styles, so that the
% pictures look alike and one convention holds throughout:
%   a FILLED endpoint belongs to the set, a HOLLOW endpoint does not.
\usepackage{tikz}
\usetikzlibrary{arrows.meta,calc,decorations.pathreplacing}
\definecolor{accent}{HTML}{A63D2F}% brick red: the set under discussion
\definecolor{calm}{HTML}{2B6A8C}%   slate blue: neighborhoods and covers
\tikzset{
  >={Stealth[length=4pt,width=3.5pt]},
  numline/.style={draw=black!55, line width=0.5pt, <->},
  tickmark/.style={draw=black!55, line width=0.5pt},
  setline/.style={draw=accent, line width=1.8pt},
  nbhdline/.style={draw=calm, line width=1.8pt},
  closedpt/.style={circle, fill=#1, draw=#1, line width=0.7pt, inner sep=0pt, minimum size=4pt},
  closedpt/.default=accent,
  openpt/.style={circle, fill=white, draw=#1, line width=0.7pt, inner sep=0pt, minimum size=4pt},
  openpt/.default=accent,
  seqpt/.style={circle, fill=black, inner sep=0pt, minimum size=2.6pt},
  lbl/.style={font=\footnotesize, inner sep=1.5pt},
  bracedeco/.style={decorate, decoration={brace, amplitude=3.5pt}, line width=0.4pt, draw=black!70},
  % graphs of functions
  graphaxis/.style={draw=black!55, line width=0.5pt, ->},
  curve/.style={draw=accent, line width=1.1pt},
  guide/.style={draw=black!45, dashed, line width=0.4pt},
  epsband/.style={fill=calm!16, draw=none},%   a tolerance on the output
  delband/.style={fill=accent!14, draw=none},% a tolerance on the input
}

% ---------------------------------------------------------------- Numbering
\setcounter{secnumdepth}{2}
\renewcommand{\thesection}{\thechapter.\arabic{section}}
\numberwithin{section}{chapter}

% \bookchapter{One}{Title} starts "Chapter One: Title".  Stepping the chapter
% counter here resets the section, theorem, definition, proposition, figure
% and equation counters automatically, so chapters need no manual \setcounter.
\newcommand{\bookchapter}[2]{%
  \stepcounter{chapter}%
  \chapter{Chapter #1: #2}}

% ---------------------------------------------------------------- Theorem environments
\newtheorem{theorem}{Theorem}[chapter]
\newtheorem{lemma}[theorem]{Lemma}
\newtheorem{corollary}[theorem]{Corollary}
\newtheorem{proposition}{Proposition}[chapter]
\theoremstyle{definition}
\newtheorem{definition}{Definition}[chapter]
\newtheorem*{remark}{Remark}
\newtheorem*{example}{Example}
\newtheorem*{imp}{Important point}
\newtheorem*{dangerousbend}{\textdbend Caution}
\newtheorem*{summary}{Summary}

% Original end-of-chapter problems: \begin{problem}[Optional title] ... \end{problem}
\newtheorem{problem}{Problem}[chapter]
% A hint, set in the margin beside a problem.
% Put it straight after \begin{problem}[...] so it lines up with the heading.
\newcommand{\hint}[1]{\leavevmode\marginnote{\emph{Hint.} #1}\ignorespaces}

% ---------------------------------------------------------------- Front-matter macros
% A full-width epigraph: \openepigraph{quotation}{author}
\newcommand{\openepigraph}[2]{%
  \begin{fullwidth}
  \sffamily\large
  \begin{doublespace}
  \noindent\allcaps{#1}\\% epigraph
  \noindent\allcaps{#2}% author
  \end{doublespace}
  \end{fullwidth}
}
.

\hypersetup{colorlinks}

% ---------------------------------------------------------------- Packages
% (The tufte-book class already loads hyperref, xcolor, paralist, titlesec,
% titletoc, multicol and natbib.)
\usepackage{amsmath, amsthm, amssymb}
\usepackage{mathrsfs}
\usepackage{booktabs}
\usepackage{enumerate}
\usepackage{manfnt}   % the "dangerous bend" sign used by the Caution environment
\usepackage{graphicx}
\setkeys{Gin}{width=\linewidth,totalheight=\textheight,keepaspectratio}
\graphicspath{{graphics/}}

\usepackage{makeidx}
\makeindex

% Cross-references to figures and tables.  Write \caption{...}\label{...} as
% usual.  The Tufte class sets a float's caption only at \end{figure}, and
% postpones the \label to match; but amsmath resets \label at the start of
% every minipage, which silently undoes that, so \ref would print the current
% section or theorem number instead of the figure number.  This restores it.
\makeatletter
\expandafter\apptocmd\csname\string\@tufte@float\endcsname
  {\let\label\@tufte@label}{}{\typeout{Warning: tufte float label patch failed}}
\makeatother

% ---------------------------------------------------------------- Figures (TikZ)
% Every diagram in the book is drawn with these shared styles, so that the
% pictures look alike and one convention holds throughout:
%   a FILLED endpoint belongs to the set, a HOLLOW endpoint does not.
\usepackage{tikz}
\usetikzlibrary{arrows.meta,calc,decorations.pathreplacing}
\definecolor{accent}{HTML}{A63D2F}% brick red: the set under discussion
\definecolor{calm}{HTML}{2B6A8C}%   slate blue: neighborhoods and covers
\tikzset{
  >={Stealth[length=4pt,width=3.5pt]},
  numline/.style={draw=black!55, line width=0.5pt, <->},
  tickmark/.style={draw=black!55, line width=0.5pt},
  setline/.style={draw=accent, line width=1.8pt},
  nbhdline/.style={draw=calm, line width=1.8pt},
  closedpt/.style={circle, fill=#1, draw=#1, line width=0.7pt, inner sep=0pt, minimum size=4pt},
  closedpt/.default=accent,
  openpt/.style={circle, fill=white, draw=#1, line width=0.7pt, inner sep=0pt, minimum size=4pt},
  openpt/.default=accent,
  seqpt/.style={circle, fill=black, inner sep=0pt, minimum size=2.6pt},
  lbl/.style={font=\footnotesize, inner sep=1.5pt},
  bracedeco/.style={decorate, decoration={brace, amplitude=3.5pt}, line width=0.4pt, draw=black!70},
  % graphs of functions
  graphaxis/.style={draw=black!55, line width=0.5pt, ->},
  curve/.style={draw=accent, line width=1.1pt},
  guide/.style={draw=black!45, dashed, line width=0.4pt},
  epsband/.style={fill=calm!16, draw=none},%   a tolerance on the output
  delband/.style={fill=accent!14, draw=none},% a tolerance on the input
}

% ---------------------------------------------------------------- Numbering
\setcounter{secnumdepth}{2}
\renewcommand{\thesection}{\thechapter.\arabic{section}}
\numberwithin{section}{chapter}

% \bookchapter{One}{Title} starts "Chapter One: Title".  Stepping the chapter
% counter here resets the section, theorem, definition, proposition, figure
% and equation counters automatically, so chapters need no manual \setcounter.
\newcommand{\bookchapter}[2]{%
  \stepcounter{chapter}%
  \chapter{Chapter #1: #2}}

% ---------------------------------------------------------------- Theorem environments
\newtheorem{theorem}{Theorem}[chapter]
\newtheorem{lemma}[theorem]{Lemma}
\newtheorem{corollary}[theorem]{Corollary}
\newtheorem{proposition}{Proposition}[chapter]
\theoremstyle{definition}
\newtheorem{definition}{Definition}[chapter]
\newtheorem*{remark}{Remark}
\newtheorem*{example}{Example}
\newtheorem*{imp}{Important point}
\newtheorem*{dangerousbend}{\textdbend Caution}
\newtheorem*{summary}{Summary}

% Original end-of-chapter problems: \begin{problem}[Optional title] ... \end{problem}
\newtheorem{problem}{Problem}[chapter]
% A hint, set in the margin beside a problem.
% Put it straight after \begin{problem}[...] so it lines up with the heading.
\newcommand{\hint}[1]{\leavevmode\marginnote{\emph{Hint.} #1}\ignorespaces}

% ---------------------------------------------------------------- Front-matter macros
% A full-width epigraph: \openepigraph{quotation}{author}
\newcommand{\openepigraph}[2]{%
  \begin{fullwidth}
  \sffamily\large
  \begin{doublespace}
  \noindent\allcaps{#1}\\% epigraph
  \noindent\allcaps{#2}% author
  \end{doublespace}
  \end{fullwidth}
}
